\subsection*{一步一步建立直观}

\paragraph{1. 让一个单元先做一件能看清的事}

输入向量不是一张图片的含义，而是被编码成数值的特征。右图只保留两个输入：权重决定各输入的影响方向和强度，偏置移动触发门槛。先计算加权和 $z$，再经过 ReLU 得到输出，这样便能把“收集证据”与“决定是否激活”分成两个步骤。



\paragraph{2. 开关边界怎样切分输入空间}

关闭单元时，加权和落在零以下，ReLU 把输出截成零。右图固定第二输入，只画输出随第一输入变化的曲线；折点对应 $z=0$，两侧分别是关闭段和激活段。拖动输入点，观察跨过折点后的变化：单个单元只产生简单的分段线性结构，多层网络通过组合这些结构形成更复杂的表示。



\paragraph{3. 训练改变规则，推理使用规则}

给定目标输出后，损失衡量当前结果与目标的差异。调权重或偏置是在改变这条输入输出规则；只拖动输入是在用既定规则计算新的样本。二者应明确区分。此处只有一个标量目标和一个单元，不代表真实分类网络的全部架构与训练过程。


